\chapter{Frames, Noise, and Quantization}
\label{ch:noise-quantization}
\fancyhead[R]{\small\textit{Chapter 11: Frames and Noise}}

\section{Overview: Closing the Loop on Robustness}
\label{sec:ch11-overview}

Chapter~1 opened this course with a promise: redundant frames are robust
not only to erasures (fulfilled in Chapter~10) but also to
\emph{noise} and \emph{quantization}. Chapter~5 began to deliver on
this promise algebraically, establishing that the expected squared
reconstruction error under additive coefficient noise is
$\sigma^2\norm{G}_F^2$ for a dual frame $G$, and that the canonical
dual minimizes this quantity. But Chapter~5 did not answer the most
practically important question: by \emph{how much} does a redundant
frame improve on a basis?

This chapter answers that question precisely and in two distinct
settings. In the \textbf{additive noise} model, each coefficient
$c_i=\inner{x}{f_i}$ is corrupted by independent noise of variance
$\sigma^2$, and the question is how the expected squared
reconstruction error scales with the oversampling ratio $r=m/n$. In
the \textbf{scalar quantization} model, each coefficient is rounded to
the nearest multiple of a step size $\delta$, and the question is how
the deterministic worst-case reconstruction error scales with $r$.

Both answers turn out to be the same: a factor of $m/n$ improvement
over a basis, achieved by the canonical dual of any tight unit-norm
frame. The two models agree because both are controlled by the same
quantity --- the Frobenius norm $\norm{\tilde F}_F^2$ of the canonical
dual's synthesis matrix --- and we close the loop Chapter~5 opened by
making this connection explicit.

\section{Additive Noise: Precise Error Accounting}
\label{sec:additive-noise}

We work with the additive noise model introduced in Chapter~5,
Section~5.6, but now carry the calculation to its conclusion.

\begin{definition}[Additive coefficient noise model]
\label{def:noise-model}
Given a frame $\{f_i\}_{i=1}^m$ for $\F^n$ and a signal $x\in\F^n$,
the receiver observes
\[
  \tilde c_i = \inner{x}{f_i} + \varepsilon_i, \qquad i=1,\ldots,m,
\]
where $\varepsilon_1,\ldots,\varepsilon_m$ are independent, mean-zero
random variables each with variance $\sigma^2$. Given a dual frame
$\{g_i\}_{i=1}^m$ with synthesis matrix $G$, the receiver reconstructs
\[
  \hat x = \sum_{i=1}^m \tilde c_i\, g_i = G\tilde c.
\]
\end{definition}

\begin{proposition}[Expected reconstruction error]
\label{prop:expected-error}
In the additive noise model, the reconstruction error satisfies
$\hat x - x = G\varepsilon$ (where $\varepsilon=(\varepsilon_i)$),
and its expected squared norm is
\begin{equation}
\label{eq:expected-error}
  \E\norm{\hat x - x}^2 = \sigma^2 \norm{G}_F^2
  = \sigma^2 \sum_{i=1}^m \norm{g_i}^2.
\end{equation}
\end{proposition}

\begin{proof}
Since $G$ reconstructs exactly from noiseless coefficients ($GF^*=I$,
Definition~5.2.1), $\hat x - x = G(F^*x+\varepsilon) - x = x + G\varepsilon
- x = G\varepsilon$. Then
\[
  \E\norm{G\varepsilon}^2 = \E\tr((G\varepsilon)(G\varepsilon)^T)
  = \tr(G\,\E[\varepsilon\varepsilon^T]\,G^T)
  = \sigma^2\tr(GG^T)
  = \sigma^2\norm{G}_F^2,
\]
using $\E[\varepsilon\varepsilon^T]=\sigma^2 I_m$ (independent, equal
variance) and $\norm{G}_F^2=\tr(GG^T)$.
\end{proof}

\begin{remark*}
This is exactly the formula derived in Chapter~5 (Section~5.6).
Proposition~\ref{prop:expected-error} is recalled here because Chapter
11's goal is to go further: to compute $\norm{G}_F^2$ explicitly for
the canonical dual of a tight unit-norm frame, and to quantify the
resulting improvement over a basis.
\end{remark*}

\subsection{The Canonical Dual Norm for Tight Unit-Norm Frames}
\label{subsec:canonical-dual-norm}

When $\{f_i\}_{i=1}^m$ is tight with bound $A$, we know from
Chapter~5 that the canonical dual is $\tilde f_i = S^{-1}f_i =
\tfrac{1}{A}f_i$ (since $S=AI$), so
$\norm{\tilde f_i} = \tfrac{1}{A}\norm{f_i}$.
For unit-norm vectors, $\norm{\tilde f_i}=\tfrac{1}{A}=\tfrac{n}{m}$.
This gives

\begin{proposition}[Canonical dual norm for tight unit-norm frames]
\label{prop:canonical-dual-norm}
Let $\{f_i\}_{i=1}^m$ be a tight, unit-norm frame for $\F^n$. Then
$A = m/n$ (trace formula, Proposition~3.4.2(iv)), and the canonical
dual's Frobenius norm satisfies
\begin{equation}
\label{eq:canonical-dual-norm}
  \norm{\tilde F}_F^2 = \sum_{i=1}^m\norm{\tilde f_i}^2
  = m\cdot\frac{1}{A^2} = m\cdot\frac{n^2}{m^2} = \frac{n^2}{m}.
\end{equation}
\end{proposition}

\begin{proof}
Each $\tilde f_i = \tfrac{1}{A}f_i$ has norm $\tfrac{1}{A}$, so
$\sum_i\norm{\tilde f_i}^2 = m/A^2 = m\cdot(n/m)^2 = n^2/m$.
\end{proof}

Substituting into Proposition~\ref{prop:expected-error}:

\begin{theorem}[Additive noise: the redundancy gain]
\label{thm:additive-noise-gain}
For a tight, unit-norm frame with $m$ vectors in $\F^n$, the canonical
dual achieves expected squared reconstruction error
\begin{equation}
\label{eq:noise-floor}
  \E\norm{\hat x - x}^2 = \sigma^2 \cdot \frac{n^2}{m} = \sigma^2\cdot
  \frac{n}{r},
  \qquad r := \frac{m}{n} \text{ (oversampling ratio)}.
\end{equation}
For comparison, a basis ($m=n$, so $r=1$) achieves expected squared
error $\sigma^2 n$.  The tight frame therefore gives a factor-of-$r$
improvement:
\[
  \frac{\E_{\text{ONB}}\norm{\hat x-x}^2}
       {\E_{\text{frame}}\norm{\hat x-x}^2} = r = \frac{m}{n}.
\]
\end{theorem}

\begin{proof}
Combine Propositions~\ref{prop:expected-error}
and~\ref{prop:canonical-dual-norm}. For the ONB baseline: $F_{\text{ONB}}
=I_n$, canonical dual $=I_n$, $\norm{I_n}_F^2 = n$, so expected error
$= \sigma^2 n = \sigma^2 n/1$ (i.e.\ $r=1$).
\end{proof}

\begin{keyidea}
\textbf{Doubling the number of frame vectors halves the noise floor.}
For a tight unit-norm frame, the expected squared reconstruction error
under additive coefficient noise is exactly $\sigma^2 n/r$, where
$r=m/n$ is the oversampling ratio. Every doubling of $r$ gives a
$3$~dB improvement in signal-to-noise ratio --- the same gain as
averaging $r$ independent measurements, which is exactly what the
canonical dual is doing: averaging over the $m$ ``views'' of $x$
provided by the frame.
\end{keyidea}

\begin{examplebox}[title={Example 11.1: Hexagonal frame in $\R^2$ ($m=6$, $r=3$)}]
\label{ex:hexagonal-noise}
The regular hexagon ($m=6$, $n=2$) is a tight unit-norm frame with
$A=3$ and canonical dual norms $\norm{\tilde f_i}=1/3$. By
Theorem~\ref{thm:additive-noise-gain}, the expected squared error is
$\sigma^2\cdot 2^2/6 = 2\sigma^2/3$, compared to $\sigma^2\cdot 2 =
2\sigma^2$ for the standard ONB. The improvement factor is exactly
$r=3$. Equivalently: to achieve the same noise floor as the ONB at
$\sigma^2$, the hexagonal frame only needs noise standard deviation
$\sigma\sqrt{r}=\sigma\sqrt{3}$ --- a $73\%$ looser noise tolerance.
\end{examplebox}

\section{The Dual-Norm Connection}
\label{sec:dual-norm-connection}

Chapter~5 (the remark following Theorem~5.5.1) promised that ``smaller
dual-vector norms translate directly into better noise suppression.''
Propositions~\ref{prop:expected-error} and~\ref{prop:canonical-dual-norm}
now make this precise: the expected squared reconstruction error is
$\sigma^2\sum_i\norm{g_i}^2$, and the canonical dual achieves the
minimum possible value of this sum, $n^2/m$, among all valid duals
(Theorem~5.5.1 of Chapter~5).

\begin{corollary}[Canonical dual is optimal for additive noise]
\label{cor:canonical-optimal}
Among all dual frames of a given frame $\{f_i\}_{i=1}^m$, the canonical
dual $\{\tilde f_i = S^{-1}f_i\}$ uniquely minimizes expected squared
reconstruction error under independent, equal-variance additive noise.
\end{corollary}

\begin{proof}
Immediate from Proposition~\ref{prop:expected-error} (error $=
\sigma^2\norm{G}_F^2$) and Theorem~5.5.1 ($\norm{\tilde F}_F^2 \le
\norm{G}_F^2$ for every dual $G$, with equality only when $G=\tilde F$).
\end{proof}

This corollary completes the loop opened in Chapter~5: the canonical
dual is optimal not merely in the abstract Frobenius-norm sense, but
concretely in the sense of minimizing the expected squared error a
receiver will observe when reconstructing from noisy measurements.

\section{Scalar Quantization}
\label{sec:scalar-quantization}

Additive noise models corruption by continuous random variables. A
different and equally important model is \emph{scalar quantization}:
each measured coefficient $\inner{x}{f_i}$ is rounded to the nearest
multiple of a step size $\delta$:
\[
  \tilde c_i = \mathcal{Q}_\delta(\inner{x}{f_i})
  := \delta\,\left\lfloor \frac{\inner{x}{f_i}}{\delta} + \frac12\right\rfloor.
\]
The quantization error $\varepsilon_i := \tilde c_i - \inner{x}{f_i}$
satisfies $|\varepsilon_i|\le\delta/2$ deterministically, but the
errors are neither independent nor Gaussian.

\begin{proposition}[Quantization error bound]
\label{prop:quantization-error}
For any dual frame $\{g_i\}$ and any scalar quantizer with step $\delta$,
the reconstruction error satisfies
\begin{equation}
\label{eq:quant-bound}
  \norm{\hat x - x}
  = \norm{G\varepsilon}
  \le \frac{\delta}{2}\sum_{i=1}^m \norm{g_i}
  \le \frac{\delta}{2}\sqrt{m}\,\norm{G}_F,
\end{equation}
where the first inequality uses $|\varepsilon_i|\le\delta/2$ and the
triangle inequality, and the second uses Cauchy--Schwarz on the sum
$\sum_i|\varepsilon_i|\norm{g_i}$.
\end{proposition}

\begin{proof}
$\norm{G\varepsilon}\le\sum_i|\varepsilon_i|\norm{g_i}\le\tfrac{\delta}{2}
\sum_i\norm{g_i}\le\tfrac{\delta}{2}\sqrt{m}\bigl(\sum_i\norm{g_i}^2\bigr)^{1/2}
=\tfrac{\delta}{2}\sqrt{m}\norm{G}_F$, by Cauchy--Schwarz.
\end{proof}

\subsection{Tight Unit-Norm Frames: The Error Bound}
\label{subsec:tight-quant-bound}

\begin{proposition}[Quantization error bound for tight unit-norm frames]
\label{prop:tight-quant-squared}
For the canonical dual of a tight unit-norm frame with $m$ vectors in
$\F^n$ and step size $\delta$, the worst-case squared reconstruction
error satisfies
\begin{equation}
\label{eq:tight-quant-bound}
  \norm{\hat x - x}^2
  = \norm{G\varepsilon}^2
  \le \norm{\varepsilon}^2\,\norm{\tilde F}_F^2
  \le m\!\left(\frac{\delta}{2}\right)^{\!2}\!\frac{n^2}{m}
  = \frac{n^2\delta^2}{4}.
\end{equation}
For a basis ($m=n$, $G=I_n$, $\norm{\varepsilon}^2 \le n(\delta/2)^2$),
the corresponding bound is also $n(\delta/2)^2 \cdot n = n^2\delta^2/4$.
Both bounds are equal: no factor-of-$m/n$ improvement is guaranteed
by the worst-case bound alone.
\end{proposition}

\begin{proof}
By the Cauchy--Schwarz inequality for vectors,
\[
  \norm{G\varepsilon}^2
  = \norm{\sum_{i=1}^m\varepsilon_i\tilde f_i}^2
  \le \norm{\varepsilon}^2\sum_{i=1}^m\norm{\tilde f_i}^2
  = \norm{\varepsilon}^2\,\norm{\tilde F}_F^2,
\]
using $\langle\sum_i a_i u_i, \sum_j a_j u_j\rangle
\le \bigl(\sum_i|a_i|^2\bigr)\bigl(\sum_i\norm{u_i}^2\bigr)$,
which follows from Cauchy--Schwarz on $\R^m$.
Since $|\varepsilon_i|\le\delta/2$ for all $i$, we have
$\norm{\varepsilon}^2\le m(\delta/2)^2$. Substituting
$\norm{\tilde F}_F^2=n^2/m$
(Proposition~\ref{prop:canonical-dual-norm}) gives the bound $n^2\delta^2/4$.
The ONB bound follows identically: $m=n$, $G=I_n$,
$\norm{I_n}_F^2=n$, $\norm\varepsilon^2\le n(\delta/2)^2$, product $n^2\delta^2/4$.
\end{proof}

\begin{remark*}
\label{rmk:quant-vs-noise}
The worst-case bound $n^2\delta^2/4$ does \emph{not} improve with $m$.
This contrast with the additive-noise result requires explanation.
In the additive noise model, each $\varepsilon_i$ is an independent
random variable with $\E[\varepsilon_i^2]=\sigma^2$; the $m$ error
contributions $\varepsilon_i\tilde f_i$ partially cancel in expectation,
giving the average squared error $\sigma^2\sum_i\norm{\tilde f_i}^2=\sigma^2 n^2/m$
--- a clear $1/m$ gain. In the quantization model with a \emph{fixed}
signal $x$, the errors $\varepsilon_i = \mathcal{Q}_\delta(\inner{x}{f_i})
- \inner{x}{f_i}$ are deterministic and can --- in the worst case --- be
systematically aligned to constructively reinforce rather than cancel,
eliminating the averaging effect. In practice, for generic signals $x$,
quantization errors do tend to behave more like independent noise and the
observed error is often close to $n^2\delta^2/(4m)$ (a factor-of-$m$
smaller); but this is an empirical observation about typical behavior,
not a guaranteed bound. The fundamental difference between average and
worst-case guarantees is a recurring theme in information theory and
signal processing.
\end{remark*}

\begin{keyidea}
The additive-noise floor $\sigma^2 n^2/m$ has clear $1/m$ scaling:
every extra frame vector reduces expected squared error by $1/m$.
The quantization bound $n^2\delta^2/4$ does \emph{not} improve with
$m$ in the worst case --- but in practice, quantization errors
behave more like noise, and empirical errors are often close to
$n^2\delta^2/(4m)$. \textbf{Both models are controlled by the same
quantity, $\norm{\tilde F}_F^2 = n^2/m$,} but the additive-noise
gain is a guaranteed bound while the quantization gain is an
empirical pattern. This distinction between average and worst-case
guarantees is fundamental in information theory.
\end{keyidea}

\section{Comparing Additive Noise and Quantization}
\label{sec:comparison}

Table~\ref{tab:comparison} summarizes the two models for a tight
unit-norm frame with $m$ vectors in $\F^n$ and its canonical dual.

\begin{table}[h]
\centering
\begin{tabular}{lcc}
\toprule
 & \textbf{Additive noise} & \textbf{Scalar quantization} \\
\midrule
Error measure & $\E\norm{\hat x-x}^2$ & $\norm{\hat x-x}^2$ (worst-case bound) \\
Controlling quantity & $\sigma^2\norm{\tilde F}_F^2$ & $m(\delta/2)^2\norm{\tilde F}_F^2$ \\
For tight unit-norm frame & $\sigma^2 n^2/m$ & $n^2\delta^2/4$ (w.c.); $\approx n^2\delta^2/(4m)$ (typ.) \\
For ONB ($m=n$) & $\sigma^2 n$ & $n^2\delta^2/4$ \\
Improvement over ONB & $m/n$ (guaranteed) & $1\times$ (w.c. bound); $m/n\times$ (empirical) \\
Minimized by & Canonical dual & Canonical dual \\
\bottomrule
\end{tabular}
\caption{Additive noise vs.\ scalar quantization for tight unit-norm frames.
The additive-noise gain is a proven average-case result; the quantization
improvement of $\approx m/n$ is empirical (worst-case bound does not depend on $m$).}
\label{tab:comparison}
\end{table}

\begin{remark*}
The additive-noise gain is $m/n$ in the expected squared error because
the $m$ noise contributions $\varepsilon_i\tilde f_i$ add incoherently
(zero-mean, independent), giving expected squared error
$\sigma^2\sum_i\norm{\tilde f_i}^2=\sigma^2 n^2/m$ --- a clear $1/m$
gain. In the quantization model, the $m$ error terms are
\emph{deterministic} and can in principle add constructively in the
worst case, eliminating the averaging effect. The proved worst-case
bound $n^2\delta^2/4$ is the same for a frame as for a basis; only
empirically (for generic signals) does the frame provide the same
$\approx m/n$ improvement as in the noise model.
\end{remark*}

\section{NumPy Implementation}
\label{sec:numpy-ch11}

\begin{lstlisting}[caption={Noise and quantization simulations for tight frames}]
import numpy as np

def tight_unit_norm_frame(n, m, seed=0):
    """
    Build a tight unit-norm frame of m vectors in R^n via gradient
    descent on the frame potential (Chapter 8).
    """
    rng = np.random.default_rng(seed)
    V = rng.standard_normal((n, m))
    V = V / np.linalg.norm(V, axis=0)
    for _ in range(2000):
        S = V @ V.T
        V = V - 0.03 * (4 * S @ V)
        V = V / np.linalg.norm(V, axis=0)
    return V

def regular_polygon(m):
    """Tight unit-norm frame in R^2: regular m-gon (m >= 3)."""
    angles = 2 * np.pi * np.arange(m) / m
    return np.vstack([np.cos(angles), np.sin(angles)])

def canonical_dual(F):
    """Canonical dual synthesis matrix: S^{-1} F."""
    return np.linalg.solve(F @ F.T, F)

def scalar_quantize(c, delta):
    """Round each entry of c to the nearest multiple of delta."""
    return delta * np.round(c / delta)

# -------------------------------------------------------------------
# Section 11.2: Additive noise -- verifying the m/n gain factor
# -------------------------------------------------------------------
sigma = 0.3
n = 2

print("=== Additive noise: expected squared error vs m/n ===")
print(f"{'m':>4}  {'r=m/n':>6}  {'predicted':>12}  {'MC mean':>10}  {'gain':>8}")
for m in [2, 3, 4, 6, 8, 12]:
    if m == 2:
        F = np.eye(2)          # ONB: the m=2 "regular polygon" is degenerate
    else:
        F = regular_polygon(m)
    G = canonical_dual(F)
    predicted = sigma**2 * np.trace(G @ G.T)

    # Monte Carlo
    rng = np.random.default_rng(42)
    x_test = np.array([1.0, -0.5])
    errs = []
    for _ in range(100_000):
        eps = rng.normal(0, sigma, m)
        x_hat = G @ (F.T @ x_test + eps)
        errs.append(np.sum((x_hat - x_test)**2))

    mc_mean = np.mean(errs)
    gain = (sigma**2 * n) / mc_mean if m > 2 else 1.0
    print(f"{m:>4}  {m/n:>6.1f}  {predicted:>12.5f}  {mc_mean:>10.5f}  {gain:>8.3f}x")

# -------------------------------------------------------------------
# Section 11.4: Scalar quantization -- verifying the proved bound n^2*delta^2/4
# and comparing to the empirically observed (approx) n^2*delta^2/(4m) behavior
# -------------------------------------------------------------------
delta = 0.2

print("\n=== Scalar quantization: proved bound vs. empirical ===")
print(f"{'m':>4}  {'proved n^2d^2/4':>16}  {'MC max':>10}  {'holds':>6}  {'n^2d^2/4m (typical)':>20}")
rng = np.random.default_rng(1)
n = 2
for m in [3, 4, 6, 8, 12]:
    F = regular_polygon(m)
    G = canonical_dual(F)
    bound = (n**2 * delta**2) / 4        # CORRECT proved bound (Prop. 11.4.1)
    typical = (n**2 * delta**2) / (4*m)  # empirical approximation only

    # Sample many x and quantize
    X = rng.standard_normal((n, 50_000))
    X = X / np.linalg.norm(X, axis=0)   # unit-norm test signals
    C_exact  = F.T @ X
    C_quant  = scalar_quantize(C_exact, delta)
    errors   = np.sum((G @ C_quant - X)**2, axis=0)
    mc_max   = errors.max()

    print(f"{m:>4}  {bound:>16.6f}  {mc_max:>10.6f}  {str(mc_max <= bound + 1e-9):>6}  {typical:>20.6f}")

# -------------------------------------------------------------------
# Example 11.1: Hexagonal frame (m=6, n=2)
# -------------------------------------------------------------------
print("\n=== Example 11.1: Hexagonal frame ===")
m, n = 6, 2
F_hex = regular_polygon(m)
G_hex = canonical_dual(F_hex)
S_hex = F_hex @ F_hex.T
A_hex = np.linalg.eigvalsh(S_hex).mean()
dual_norms = np.linalg.norm(G_hex, axis=0)
print(f"  Frame bound A = {A_hex:.4f}  (expected m/n = {m/n:.4f})")
print(f"  Dual-vector norms: {dual_norms.round(4)}  (expected n/m = {n/m:.4f})")
print(f"  ||tilde_F||_F^2 = {np.trace(G_hex @ G_hex.T):.4f}"
      f"  (expected n^2/m = {n**2/m:.4f})")
sigma_test = 0.3
print(f"  Predicted noise floor: sigma^2 * n^2/m ="
      f" {sigma_test**2 * n**2/m:.4f}")
print(f"  ONB comparison:        sigma^2 * n   ="
      f" {sigma_test**2 * n:.4f}")
print(f"  Improvement factor: {(sigma_test**2 * n) / (sigma_test**2 * n**2/m):.1f}x"
      f"  (= m/n = {m/n:.1f})")
\end{lstlisting}

\section{Chapter Summary}
\label{sec:summary-ch11}

\begin{itemize}
  \item \textbf{Additive noise model.} The expected squared reconstruction
        error using dual frame $G$ is $\sigma^2\norm{G}_F^2$
        (Proposition~\ref{prop:expected-error}). The canonical dual
        uniquely minimizes this (Corollary~\ref{cor:canonical-optimal},
        closing the loop from Chapter~5).

  \item \textbf{Tight unit-norm frames.} The canonical dual has
        $\norm{\tilde F}_F^2 = n^2/m$
        (Proposition~\ref{prop:canonical-dual-norm}), giving expected
        squared error $\sigma^2 n^2/m = \sigma^2 n/r$ where $r=m/n$
        (Theorem~\ref{thm:additive-noise-gain}). This is a factor-of-$r$
        improvement over an ONB.

  \item \textbf{Scalar quantization.} With step $\delta$ and the
        canonical dual, the proved worst-case squared reconstruction
        error is bounded by $n^2\delta^2/4$
        (Proposition~\ref{prop:tight-quant-squared}) --- the same as for
        a basis. The factor-of-$m/n$ improvement is observed empirically
        (typical errors $\approx n^2\delta^2/(4m)$) but is not guaranteed
        in the worst case; see the remark in Section~\ref{subsec:tight-quant-bound}.

  \item \textbf{Both models, same mechanism.} The factor-of-$m/n$
        improvement comes from the same source in both models: dual
        vectors of norm $n/m$ (rather than norm $1$ for a basis), so each
        error contribution $\varepsilon_i\tilde f_i$ is smaller by $n/m$,
        and there are $m$ of them. The controlling quantity is
        $\norm{\tilde F}_F^2 = n^2/m$ in both cases (Table~\ref{tab:comparison}).
\end{itemize}

\section{Lab Exercises}
\label{sec:lab-ch11}

\begin{exercisebox}[title={Lab Exercise 11.1: Reproducing the Noise Floor Table}]
Implement \texttt{regular\_polygon}, \texttt{canonical\_dual}, and the
Monte Carlo loop from Section~\ref{sec:numpy-ch11}. For $n=2$ and
$m=2,3,4,6,8,12,20$, compute and tabulate: (a)~the predicted noise
floor $\sigma^2 n^2/m$, (b)~the Monte Carlo mean squared error (use
$\sigma=0.3$, $10^5$ trials), and (c)~the gain factor over the ONB.
Verify that the gain equals $m/n$ to at least two decimal places for
all $m\ge3$.
\end{exercisebox}

\begin{exercisebox}[title={Lab Exercise 11.2: Non-Tight Frame Penalty}]
Consider the non-tight frame $\{h_1,h_2,h_3\}$ with $h_1=e_1$,
$h_2=e_2$, $h_3=(e_1+e_2)/\sqrt{2}$ in $\R^2$ (frame bounds $A=1$,
$B=2$; note this differs from Chapter~5's frame by normalizing $h_3$):
\begin{enumerate}[label=(\alph*)]
  \item Compute the canonical dual $\tilde H$ and its Frobenius norm
        $\norm{\tilde H}_F^2$.
  \item Compare $\norm{\tilde H}_F^2$ to the tight unit-norm prediction
        $n^2/m = 4/3$ for $n=2$, $m=3$. Which is larger? Does the
        non-tight frame get a better or worse noise floor than the
        tight triangular frame with the same $m$?
  \item Run $10^5$ Monte Carlo noise trials (with $\sigma=0.3$) and
        verify your prediction from~(a).
  \item Explain why the non-tight frame performs worse: what is it
        about the frame geometry (not just the formula) that leads to
        a larger dual Frobenius norm?
\end{enumerate}
\end{exercisebox}

\begin{exercisebox}[title={Lab Exercise 11.3: Quantization Simulation}]
For $n=2$ and $m \in \{3,4,6,8,12\}$, use the code in
Section~\ref{sec:numpy-ch11} to:
\begin{enumerate}[label=(\alph*)]
  \item Verify that the \emph{proved} worst-case quantization bound
        $n^2\delta^2/4$ (Proposition~\ref{prop:tight-quant-squared}) is
        never violated across $10^5$ random unit-norm test signals
        (use $\delta=0.2$). Confirm this holds for all $m$.
  \item Plot the empirical maximum error (over all test signals) as a
        function of $m$ on the same axes as the proved bound
        $n^2\delta^2/4$ (constant in $m$) and the empirical approximation
        $n^2\delta^2/(4m)$ (decays as $1/m$). Which does the empirical
        maximum track more closely?
  \item The proved bound is a \emph{constant} (independent of $m$), but
        the empirical maximum decays with $m$. Explain why in terms of the
        remark in Section~\ref{subsec:tight-quant-bound}: when does the
        worst-case alignment of quantization errors actually arise?
        Then try to construct, for a fixed frame ($m=3$, $n=2$), a signal
        $x$ and step size $\delta$ such that the quantization errors
        $\varepsilon_i$ all have the same sign and the reconstruction error
        is noticeably larger than $n^2\delta^2/(4m)$. Report the ratio
        of actual to ``typical'' error for your constructed $x$.
\end{enumerate}
\end{exercisebox}

\begin{exercisebox}[title={Lab Exercise 11.4: Higher Dimensions}]
Extend the analysis to $n=3$ and $m=6,9,12,18$. Use
\texttt{tight\_unit\_norm\_frame} (gradient-descent construction from
Section~\ref{sec:numpy-ch11}) to build tight frames for $m>5$ where
regular-polygon constructions no longer apply. For each $(n,m)$ pair:
\begin{enumerate}[label=(\alph*)]
  \item Verify $\norm{\tilde F}_F^2 \approx n^2/m$ for the gradient-descent
        tight frame.
  \item Run Monte Carlo noise trials and compare to the predicted noise
        floor.
  \item Plot the noise floor as a function of $m$ for fixed $n=3$, and
        compare the slope to the theoretical $\sim 1/m$ decay.
\end{enumerate}
\end{exercisebox}

\begin{exercisebox}[title={Lab Exercise 11.5 (Challenge): Alternate Duals and Noise}]
Theorem~5.5.1 of Chapter~5 guarantees the canonical dual minimizes
$\norm{G}_F^2$ and hence the noise floor. How much worse can alternate
duals be?
\begin{enumerate}[label=(\alph*)]
  \item For the triangular frame, parametrize the full dual family
        $G(w) = \tilde F + wk^T$ (Chapter~6, Proposition~6.2.1) and
        compute $\norm{G(w)}_F^2 = 4/3 + \norm{w}^2$ as a function of
        $w\in\R^2$.
  \item For $\norm{w}=1$ (one unit away from the canonical dual in the
        parameter space), compute the predicted noise floor
        $\sigma^2(4/3+1)$ and compare to the canonical minimum
        $\sigma^2\cdot 4/3$. What is the percentage penalty for using
        this sub-optimal dual?
  \item The Chapter~6 analysis showed there is a special point
        $w^*=(1,1)/\sqrt3$ where the error ellipse for the $\{h_1,h_2,h_3\}$
        frame becomes circular (isotropic). Compute $\norm{G(w^*)}_F^2$
        for this point and compare to the canonical dual's noise floor.
        This illustrates that minimizing total noise and minimizing noise
        anisotropy are two different objectives that generally require
        two different dual choices.
\end{enumerate}
\end{exercisebox}
\newpage
\section{Supplementary Exercises}
\label{sec:extra-ch11}

\subsection*{Theoretical Exercises}
\begin{theoexercises}
	\item Suppose noise is added directly to the \emph{signal} $x$ (i.e.\
	the receiver observes $x+\varepsilon$ for $\varepsilon\in\R^n$,
	\emph{before} any frame analysis takes place), rather than to the
	frame coefficients. Prove that in this model, redundancy provides
	\emph{no} benefit whatsoever: the reconstruction error is
	$\norm{\varepsilon}$ regardless of which frame or dual is used.
	Contrast with Theorem~11.2.4, and identify precisely where in the
	chain $x\to(\inner{x}{f_i})_i\to\hat x$ redundancy's benefit
	actually enters.
	
	\item Prove that the \emph{worst-case} quantization error is controlled
	by the operator norm $\norm{G}_{\mathrm{op}}$ (largest singular
	value) of the dual synthesis matrix, not the Frobenius norm
	$\norm{G}_F$ used for \emph{expected} additive-noise error: show
	$\norm{G\varepsilon}_2\le\norm{G}_{\mathrm{op}}\norm{\varepsilon}_2$
	for any bounded quantization error vector $\varepsilon$.
	
	\item Prove that for a tight, unit-norm frame with $m$ vectors in
	$\R^n$, the canonical dual's operator norm is exactly
	$\norm{\tilde F}_{\mathrm{op}}=\sqrt{n/m}$. Compare the rate at
	which this quantity decays in $m$ (like $1/\sqrt m$) to the rate
	at which the expected-error-controlling quantity
	$\norm{\tilde F}_F^2=n^2/m$ decays (like $1/m$), and explain the
	practical implication for worst-case versus average-case
	robustness.
	
	\item Prove the general (not necessarily tight) two-sided bound
	$1/\sqrt B\le\norm{\tilde F}_{\mathrm{op}}\le1/\sqrt A$ for the
	canonical dual's operator norm, in terms of the original frame's
	bounds $A,B$.
	
	\item For the $m=n+1$ dual family $G(w)=\tilde F+wk^T$ of Chapter~6,
	prove or disprove: the operator norm $\norm{G(w)}_{\mathrm{op}}$,
	like the Frobenius norm, is uniquely minimized at $w=0$.
\end{theoexercises}

\subsection*{Computational Exercises}
\begin{compexercises}
	\item For $n=2$, $m=3,6,12,24$, compute $\norm{\tilde F}_{\mathrm{op}}$
	for the regular $m$-gon and confirm it matches $\sqrt{n/m}$
	(T3), plotting both $\norm{\tilde F}_{\mathrm{op}}$ and
	$\norm{\tilde F}_F^2$ against $m$ on a log-log scale to compare
	decay rates.
	
	\item For $\{h_1,h_2,h_3\}$, compute $\norm{\tilde F}_{\mathrm{op}}$
	directly and confirm it lies within the bounds $1/\sqrt B$ and
	$1/\sqrt A$ from T4.
	
	\item For the triangular frame's dual family $G(w)$, compute
	$\norm{G(w)}_{\mathrm{op}}$ over a grid of $w$ values and confirm
	numerically whether it is minimized at $w=0$, resolving T5.
	
	\item Simulate the signal-level noise model from T1 directly: add
	Gaussian noise to $x$ itself, ``reconstruct'' by simply returning
	the noisy observation, and confirm the error equals
	$\norm{\varepsilon}$ regardless of which frame's coefficients (if
	any) are computed downstream.
\end{compexercises}